\documentclass[14pt]{beamer}
\usepackage{pepslides}
\title{Finite elasticity}
\lecturenumber{13}
\hypersetup{pdftitle={Lecture 13: Finite elasticity}}

\begin{document}

\titleslide

\begin{frame}
\centering
\includegraphics[height=0.85\paperheight]{tromp2010_fig7_body_waves.png}
\credit{Tromp \emph{et al.} (2010)}
\end{frame}

\begin{frame}
\centering
\begin{tabular}{cc}
\includegraphics[height=0.42\paperheight]{davies2022_convection_pacific.png} &
\includegraphics[width=0.47\textwidth]{seismicity_M5_1976-2025.png} \\[1em]
\includegraphics[width=0.47\textwidth]{pugmire_glad-m25.jpg} &
\includegraphics[width=0.47\textwidth]{barbot2023_turkiye_slip.png}
\end{tabular}
\credit{Davies \emph{et al.} (2022); Pugmire (2025); USGS (2025); Barbot \emph{et al.} (2023)}
\end{frame}

\sectionslide{Kinematics}

\begin{frame}
\centering
\includegraphics[height=0.5\paperheight]{L1F1.pdf}

\raggedright
Deformation of a continuum defined relative to a reference body.

The \textbf{motion} is a field that specifies the position of particle $\mathbf{x}$ at time $t$ in physical space.

The motion is assumed to be smooth and have a smooth inverse -- not fracturing or tearing or interpenetration.

The conventional choice of reference body is based on a reference time:
\begin{eqn}
\bphi(\mathbf{x},t_{0})=\mathbf{x}.
\end{eqn}
\end{frame}

\begin{frame}
The (referential) \textbf{velocity} is the time derivative of the motion:
\begin{eqn}
v_{i}(\mathbf{x},t)=\frac{\partial\varphi_{i}}{\partial t}(\mathbf{x},t).
\end{eqn}
The \textbf{deformation gradient} is the spatial derivative of the motion:
\begin{eqn}
\begin{gathered}
F_{ij}=\frac{\partial\varphi_{i}}{\partial x_{j}},\\[\eqgap]
\varphi_{i}(\mathbf{x}+\delta\mathbf{x},t)=\varphi_{i}(\mathbf{x},t)+F_{ij}(\mathbf{x},t)\delta x_{j}+\cdots.
\end{gathered}
\end{eqn}
Because the motion maps the reference body one-to-one with the instantaneous body, the deformation gradient
is at each point an invertible matrix. We are free to assume its determinant is \textbf{everywhere positive}.
\end{frame}

\begin{frame}
The \textbf{Jacobian} of the motion is defined by:
\begin{eqn}
J(\mathbf{x},t)=\det\mathbf{F}(\mathbf{x},t).
\end{eqn}
We can use it to change variables within volume integrals:
\begin{eqn}
\begin{gathered}
I=\int_{M_{t}}f(\mathbf{y})\dd^{3}\mathbf{y},\\[\eqgap]
I=\int_{M}f[\bphi(\mathbf{x},t)]J(\mathbf{x},t)\dd^{3}\mathbf{x}.
\end{gathered}
\end{eqn}
In this manner, we remove the need to explicitly consider moving boundaries.
\end{frame}

\begin{frame}
Mass of a sub-body does not change with time:
\begin{eqn}
\frac{\ddns}{\ddns t}\int_{U_{t}}\varrho(\mathbf{y},t)\dd^{3}\mathbf{y}=0.
\end{eqn}
Transforming this to the reference body, we get:
\begin{eqn}
\frac{\ddns}{\ddns t}\int_{U}J(\mathbf{x},t)\varrho[\bphi(\mathbf{x},t),t]\dd^{3}\mathbf{x}=0.
\end{eqn}
The reference volume is arbitrary, and so the integrand does not depend on time:
\begin{eqn}
\rho(\mathbf{x})=J(\mathbf{x},t)\varrho[\bphi(\mathbf{x},t),t].
\end{eqn}
The \textbf{referential density} equals the density at the reference time.
\end{frame}

\sectionslide{Dynamics}

\begin{frame}
Kinetic energy density per unit volume in the reference body:
\begin{eqn}
\frac{1}{2}\rho(\mathbf{x})\|\mathbf{v}(\mathbf{x},t)\|^{2}.
\end{eqn}
Total kinetic energy associated with a motion:
\begin{eqn}
\mathcal{T}(t)=\frac{1}{2}\int_{M}\rho(\mathbf{x})\|\mathbf{v}(\mathbf{x},t)\|^{2}\dd^{3}\mathbf{x}.
\end{eqn}
This is just generalising particle mechanics to a continuum. What about the potential energy?
\end{frame}

\begin{frame}
As a starting point, we assume that the potential energy can be written:
\begin{eqn}
\mathcal{V}(t)=\int_{M}W[\mathbf{x},t,\bphi(\mathbf{x},t),\mathbf{F}(\mathbf{x},t)]\dd^{3}\mathbf{x}.
\end{eqn}
This choice defines a \textbf{simple hyperelastic material}:
\begin{itemize}
\item Hyperelastic -- a potential energy for the body exists (cf.\ conservative forces in particle mechanics)
\item Simple -- the potential energy density depends only on first spatial derivatives of the motion
\end{itemize}
The key physical intuition is that elastic materials have a preferred shape or volume, and that energy is required
to move them into a different configuration.
\end{frame}

\begin{frame}
\textbf{Material frame indifference} requires the potential energy be invariant under a set of transformations:

\textbf{Time-translation:}
\begin{eqn}
\begin{gathered}
\bphi^{\prime}(\mathbf{x},t)=\bphi(\mathbf{x},t-\tau)\arrowto\mathcal{V}^{\prime}(t)=\mathcal{V}(t-\tau),\\[\eqgap]
\mathcal{V}(t)=\int_{M}W[\mathbf{x},\bphi(\mathbf{x},t),\mathbf{F}(\mathbf{x},t)]\dd^{3}\mathbf{x}.
\end{gathered}
\end{eqn}
\textbf{Rigid-body motions:}
\begin{eqn}
\begin{gathered}
\bphi^{\prime}(\mathbf{x},t)=\mathbf{a}(t)+\mathbf{Q}(t)\bphi(\mathbf{x},t)\arrowto\mathcal{V}^{\prime}(t)=\mathcal{V}(t),\\[\eqgap]
W[\mathbf{x},\mathbf{F}(\mathbf{x},t)]=W[\mathbf{x},\mathbf{Q}(t)\mathbf{F}(\mathbf{x},t)].
\end{gathered}
\end{eqn}
Invariance under superimposed rigid rotations discussed further next time.
\end{frame}

\begin{frame}
The \textbf{Lagrangian density} for a simple hyperelastic body is defined by:
\begin{eqn}
\mathcal{L}(\mathbf{x},t,\bphi,\mathbf{v},\mathbf{F})=\frac{1}{2}\rho(\mathbf{x})\|\mathbf{v}\|^{2}-W(\mathbf{x},\mathbf{F}).
\end{eqn}
The corresponding \textbf{action} is given by:
\begin{eqn}
\mathcal{S}(\bphi)=\int_{0}^{T}\int_{M}\mathcal{L}[\mathbf{x},t,\bphi(\mathbf{x},t),\mathbf{v}(\mathbf{x},t),\mathbf{F}(\mathbf{x},t)]\dd^{3}\mathbf{x}\dd t.
\end{eqn}
\textbf{Hamilton's principle} states that the true motion is a \textbf{stationary point} of the action for variations subject
to fixed end-point conditions. This means that the \textbf{first variation of the action}
\begin{eqn}
\delta\mathcal{S}=\lim_{h\rightarrow0}\frac{1}{h}[\mathcal{S}(\bphi+h\,\delta\bphi)-\mathcal{S}(\bphi)]
\end{eqn}
is zero for all perturbations to the motion such that:
\begin{eqn}
\delta\bphi(\mathbf{x},0)=\delta\bphi(\mathbf{x},T)=\mathbf{0}.
\end{eqn}
\end{frame}

\begin{frame}
From Hamilton's principle, we arrive at the \textbf{Euler--Lagrange equations}:
\begin{eqn}
\frac{\partial\mathcal{L}}{\partial\varphi_{i}}-\frac{\partial}{\partial t}\left(\frac{\partial\mathcal{L}}{\partial v_{i}}\right)-\frac{\partial}{\partial x_{j}}\left(\frac{\partial\mathcal{L}}{\partial F_{ij}}\right)=0,
\end{eqn}
along with the \textbf{natural boundary conditions}:
\begin{eqn}
\frac{\partial\mathcal{L}}{\partial F_{ij}}\hat{n}_{j}=0.
\end{eqn}
Full details of the derivation are given in the handout.
\end{frame}

\begin{frame}
Using the concrete form for the Lagrangian density, the equations of motion are:
\begin{eqn}
\rho\frac{\partial v_{i}}{\partial t}-\frac{\partial T_{ij}}{\partial x_{j}}=0,
\end{eqn}
along with the traction-free boundary conditions:
\begin{eqn}
T_{ij}\hat{n}_{j}=0,
\end{eqn}
where we have defined the \textbf{first Piola--Kirchhoff stress} by:
\begin{eqn}
T_{ij}(\mathbf{x},t)=\frac{\partial W}{\partial F_{ij}}[\mathbf{x},\mathbf{F}(\mathbf{x},t)].
\end{eqn}
This tensor is the referential counterpart of the \textbf{Cauchy stress} seen within spatial (Eulerian) formulations of
continuum mechanics.

In a referential (Lagrangian) framework, conservation of mass is trivially met, while the equations of motion are
posed on a fixed reference body.
\end{frame}

\begin{frame}
\centering
\begin{tikzpicture}[line width=0.9pt,>={Stealth[length=3mm]},scale=1.4]
  \draw plot[smooth cycle, tension=0.8] coordinates
    {(-4.4,0.4) (-3.6,1.5) (-2.0,1.3) (-1.2,0.3) (-1.6,-1.2) (-3.2,-1.5) (-4.3,-0.6)};
  \node at (-3.7,-0.8) {\large $M$};
  \draw[line width=2.2pt] (-3.0,-0.3) -- (-2.2,0.3);
  \draw[->] (-2.6,0.0) -- (-3.14,0.72) node[left] {$\hat{\mathbf{n}}$};
  \node[below right] at (-2.45,-0.1) {$\ddns S$};
  \draw plot[smooth cycle, tension=0.8] coordinates
    {(1.2,0.9) (2.2,1.9) (4.0,1.7) (5.0,0.5) (4.4,-1.3) (2.8,-1.7) (1.4,-0.8)};
  \node at (4.0,-0.9) {\large $M_{t}$};
  \draw[line width=2.2pt] (2.3,0.1) -- (3.7,0.45);
  \draw[->] (3.0,0.275) -- (2.79,1.15) node[above left] {$\hat{\mathbf{n}}_{t}$};
  \draw[->,line width=1.4pt] (3.0,0.275) -- (4.05,1.25) node[right] {$\ddns\mathbf{F}$};
  \node[below] at (3.15,0.15) {$\ddns S_{t}$};
  \draw[->] (-2.2,-0.9) to[bend right=25] node[below,midway] {$\bm{\varphi}(\cdot,t)$} (2.6,-0.9);
\end{tikzpicture}
\begin{eqn}
\ddns\mathbf{F}=\mathbf{T}\hat{\mathbf{n}}\dd S=\bm{\sigma}\hat{\mathbf{n}}_{t}\dd S_{t}.
\end{eqn}
\raggedright
The first Piola--Kirchhoff stress, $\mathbf{T}$, gives the force on the deformed surface element per unit area of the
reference element; the Cauchy stress, $\bm{\sigma}$, gives the same force per unit area of the deformed element.
\end{frame}

\begin{frame}[t]
\heading{What you need to know and be able to do}
\begin{enumerate}
\item[(i)] Understand the terminology needed to describe the kinematics of an elastic body.
\item[(ii)] Use the conservation of mass to derive the relation between the spatial and referential densities.
\item[(iii)] State expressions for the kinetic and potential energies of a simple hyperelastic body, and explain the
underlying physical ideas and assumptions.
\item[(iv)] Give a statement of Hamilton's principle in the context of finite elasticity, and use it to derive the
equations of motion and associated boundary conditions.
\end{enumerate}
\end{frame}

\end{document}
